\documentclass[11pt]{article}

\usepackage[utf8]{inputenc}
\usepackage[T1]{fontenc}
\usepackage{lmodern}
\usepackage[margin=1in]{geometry}
\usepackage{amsmath, amssymb, amsfonts, mathtools}
\usepackage{bm}
\usepackage{microtype}
\usepackage{xcolor}
\usepackage{hyperref}
\usepackage{enumitem}
\usepackage{booktabs}
\usepackage{parskip}

\hypersetup{
    colorlinks=true,
    linkcolor=blue,
    urlcolor=blue,
    citecolor=blue
}

\title{Model Brief For Presentation}
\author{}
\date{}

\begin{document}

\maketitle

\section*{Main Idea}

The core model is an audio-conditioned autoregressive transformer for osu!taiko chart generation.

\begin{itemize}[leftmargin=1.5em]
    \item Input: beat-aligned mel-spectrogram window $X \in \mathbb{R}^{192 \times 128}$
    \item Output: token sequence $\mathbf{y} = (y_1, \dots, y_T)$ containing timing tokens such as \texttt{TS\_*} and event tokens such as \texttt{DON}, \texttt{KAT}, \texttt{BIGDON}, and \texttt{BIGKAT}
    \item Goal: model the conditional distribution
\end{itemize}

\[
p(\mathbf{y} \mid X) = \prod_{t=1}^{T} p(y_t \mid y_{<t}, X)
\]

So the model listens to a short audio window, then writes the chart token by token.

\section*{AR Transformer}

The baseline is an encoder-decoder transformer.

\begin{itemize}[leftmargin=1.5em]
    \item Audio encoder:
    \begin{itemize}[leftmargin=1.5em]
        \item project mel features into model space
        \item add positional encoding
        \item encode with transformer self-attention
    \end{itemize}
    \item Token decoder:
    \begin{itemize}[leftmargin=1.5em]
        \item embed previous tokens
        \item add positional encoding
        \item apply causal self-attention
        \item cross-attend to encoded audio
        \item predict the next token
    \end{itemize}
\end{itemize}

Compact math:

\[
H = \mathrm{Enc}(\mathrm{PE}(X W_a))
\]

\[
E_{<t} = \mathrm{PE}(\mathrm{TokEmb}(y_{<t}))
\]

\[
h_t = \mathrm{Dec}(E_{<t}, H)
\]

\[
p(y_t \mid y_{<t}, X) = \mathrm{softmax}(W_o h_t + b_o)
\]

Training objective:

\[
\mathcal{L}_{\mathrm{AR}} = - \sum_{t=1}^{T} \log p(y_t^{\ast} \mid y_{<t}^{\ast}, X)
\]

This is standard teacher-forced cross-entropy over the target token sequence.

\section*{Metadata Conditioning}

The decoder also receives three global conditioning signals:

\begin{itemize}[leftmargin=1.5em]
    \item difficulty
    \item density
    \item beatmap ID as a style / identity proxy
\end{itemize}

Let:

\begin{itemize}[leftmargin=1.5em]
    \item $d_{\mathrm{diff}}$ be the difficulty number after scaling it into a range the model can use
    \item $d_{\mathrm{dens}}$ be the density number after scaling it into a range the model can use
    \item $d_{\mathrm{id}}$ be the beatmap-ID number after scaling it, used as a rough style hint
    \item $\mathbf{c}$ be the 3-number column made by putting those values together
    \item $\phi(\cdot)$ be the small neural network that turns those 3 numbers into a richer feature
    \item $\mathbf{g}$ be the output of that small neural network, meaning the final metadata feature the model uses
    \item $E_{<t}$ be the decoder's token representation before metadata is added
    \item $\tilde{E}_{<t}$ be the decoder's token representation after metadata is added
    \item $\alpha$ be a learned number that controls how strongly the metadata affects the decoder
\end{itemize}

These are normalized, stacked, projected, and added to decoder token states:

\[
\mathbf{c} = [d_{\mathrm{diff}}, d_{\mathrm{dens}}, d_{\mathrm{id}}]^{\top}
\]

\[
\mathbf{g} = \phi(\mathbf{c})
\]

\[
\tilde{E}_{<t} = E_{<t} + \alpha \mathbf{g}
\]

Interpretation of the symbols:

\begin{itemize}[leftmargin=1.5em]
    \item $\mathbf{c} \in \mathbb{R}^{3}$ just means ``a list of 3 real numbers''
    \item $\mathbf{g}$ is the model's learned summary of the metadata
    \item the same $\mathbf{g}$ is copied to every decoder step, because the whole chart window shares the same metadata
    \item $\alpha$ lets the model decide whether metadata should matter a little or a lot
\end{itemize}

Interpretation:

\begin{itemize}[leftmargin=1.5em]
    \item difficulty biases chart hardness
    \item density biases note rate
    \item beatmap ID acts as a weak style prior for mapper/chart tendencies
\end{itemize}

\section*{Context Transformer Variation}

The context model keeps the same AR backbone, but changes the decoder input context.

Instead of decoding from only the current window, it prepends:

\begin{itemize}[leftmargin=1.5em]
    \item recent history tokens
    \item retrieved motif tokens from earlier similar windows
    \item current-window tokens
\end{itemize}

Segment embeddings tell the decoder which tokens come from which source.

This is meant to improve:

\begin{itemize}[leftmargin=1.5em]
    \item cross-window continuity
    \item motif recurrence
    \item phrase-level consistency
\end{itemize}

\section*{Retrieval Math}

Define the symbols locally:

\begin{itemize}[leftmargin=1.5em]
    \item $H_i$ is the encoder output for an earlier audio window $i$
    \item $H_{\mathrm{cur}}$ is the encoder output for the current audio window the model is working on now
    \item $\mathrm{mean}(H_i)$ means ``average all the time-step vectors in window $i$ into one summary vector''
    \item $\mathbf{k}_i$ is the stored lookup vector for earlier window $i$
    \item $\mathbf{q}$ is the lookup vector for the current window
    \item $s_i$ is the similarity score between the current window and earlier window $i$
    \item $K$ is how many earlier windows we keep after retrieval
    \item $\mathbf{z}_{\mathrm{hist}}$ is the block of token representations from recent chart history
    \item $\mathbf{z}_{\mathrm{retr}}$ is the block of token representations copied from retrieved earlier motifs
    \item $\mathbf{z}_{\mathrm{cur}}$ is the block of token representations for the current part being decoded
    \item $\mathbf{z}_{\mathrm{ctx}}$ is the full decoder input after joining history, retrieval, and current tokens together
\end{itemize}

For each prior window, the model computes a pooled audio key from the encoded memory:

\[
\mathbf{k}_i = \frac{\mathrm{mean}(H_i)}{\lVert \mathrm{mean}(H_i) \rVert_2}
\]

For the current window:

\[
\mathbf{q} = \frac{\mathrm{mean}(H_{\mathrm{cur}})}{\lVert \mathrm{mean}(H_{\mathrm{cur}}) \rVert_2}
\]

Similarity is cosine similarity:

\[
s_i = \mathbf{q}^{\top} \mathbf{k}_i
\]

Interpretation of the symbols:

\begin{itemize}[leftmargin=1.5em]
    \item $\mathbf{k}_i$ is the ``saved fingerprint'' for earlier window $i$
    \item $\mathbf{q}$ is the ``current fingerprint'' used to search for similar earlier windows
    \item $s_i$ gets larger when the current window and earlier window look more alike
    \item this is cosine similarity because both vectors are scaled to length 1 before comparison
\end{itemize}

Then retrieve the top $K$ prior windows with highest $s_i$, excluding the most recent few windows to avoid trivial copying.

The decoder input becomes:

\[
\mathbf{z}_{\mathrm{ctx}} = [\mathbf{z}_{\mathrm{hist}} ; \mathbf{z}_{\mathrm{retr}} ; \mathbf{z}_{\mathrm{cur}}]
\]

Here $[\cdot ; \cdot ; \cdot]$ denotes concatenation along the token dimension.

In simpler words, this means the decoder reads:

\begin{itemize}[leftmargin=1.5em]
    \item some recent chart history
    \item some retrieved old patterns that match the current music
    \item the current generation region
\end{itemize}

all as one long input sequence.

So the context model is effectively using retrieval-augmented sequence modeling over earlier musical structure.

\section*{Key Optimizations}

\subsection*{Data / Training Setup}

\begin{itemize}[leftmargin=1.5em]
    \item Beat-aligned representation at \texttt{48} frames per beat
    \begin{itemize}[leftmargin=1.5em]
        \item This turns raw audio into a rhythm-aware grid, so the model sees music in units that line up with beats instead of arbitrary waveform chunks.
        \item That makes the token timing task much easier, because note placement is learned relative to beat structure.
    \end{itemize}
    \item 4-beat training windows
    \begin{itemize}[leftmargin=1.5em]
        \item The repo uses short fixed windows with shape $(192, 128)$.
        \item This keeps the problem computationally manageable and gives stable fixed-size supervision for the AR model.
    \end{itemize}
    \item constant-BPM taiko-only filtering
    \begin{itemize}[leftmargin=1.5em]
        \item Charts with BPM changes are removed from the main training/inference path.
        \item This simplifies beat alignment, timing-token generation, and reconstruction, reducing noise in the training data.
    \end{itemize}
    \item optional \texttt{keep\_only\_max\_notes\_per\_song} curation
    \begin{itemize}[leftmargin=1.5em]
        \item When multiple charts exist for the same song, this option keeps only the densest chart.
        \item The idea is to reduce conflicting supervision from many alternate difficulties and focus training on one strong target per song.
    \end{itemize}
    \item snapshot datasets for faster iteration
    \begin{itemize}[leftmargin=1.5em]
        \item Instead of always training on the full eligible pool, the repo can build a curated snapshot subset.
        \item This makes experiments cheaper and faster while still using the same preprocessing and model code.
    \end{itemize}
\end{itemize}

\subsection*{Runtime / Engineering}

\begin{itemize}[leftmargin=1.5em]
    \item cached manifests, splits, and sequence indexes
    \begin{itemize}[leftmargin=1.5em]
        \item The repo caches chart manifests and train/validation/test sequence indexes.
        \item This avoids repeating expensive bookkeeping every time a new run starts, which speeds up experimentation and improves reproducibility.
    \end{itemize}
    \item mixed precision (\texttt{bf16} / \texttt{fp16} when available)
    \begin{itemize}[leftmargin=1.5em]
        \item Training can automatically use lower-precision math on supported hardware.
        \item This reduces memory use and often increases throughput, allowing larger batches or faster epochs without changing the model design.
    \end{itemize}
    \item tuned dataloader settings
    \begin{itemize}[leftmargin=1.5em]
        \item Settings such as \texttt{pin\_memory}, worker count, prefetching, and persistent workers are exposed and used in optimized runs.
        \item These are not model changes, but they matter because slow data loading can bottleneck GPU training.
    \end{itemize}
    \item exported inference snapshots during training
    \begin{itemize}[leftmargin=1.5em]
        \item The training loop can periodically save inference-ready checkpoints instead of only saving the final model.
        \item This lets the team test intermediate models directly in the real inference path and compare quality before the end of training.
    \end{itemize}
    \item throughput and adherence monitoring
    \begin{itemize}[leftmargin=1.5em]
        \item The repo tracks runtime measures such as tokens-per-second as well as proxy adherence metrics such as density error and difficulty drift.
        \item These diagnostics help evaluate not just whether loss is decreasing, but whether the conditioning behavior is staying aligned with the intended chart properties.
    \end{itemize}
\end{itemize}

\subsection*{Why These Optimizations Matter}

These optimizations matter because the project is not only about designing a model, but also about making the full training-and-inference workflow practical.

\begin{itemize}[leftmargin=1.5em]
    \item Data filtering reduces label noise and keeps the task well-defined.
    \item Windowing and beat alignment make the learning problem simpler.
    \item Snapshot datasets and maxopt curation make experiments faster and more controlled.
    \item Mixed precision and dataloader tuning improve hardware efficiency.
    \item Cached indexes and inference snapshots reduce iteration cost.
\end{itemize}

So the optimization story is: simplify the learning problem, clean the supervision, and make iteration fast enough that many model variants can actually be tested.

\section*{One-Line Takeaway}

The baseline model is an audio-to-token autoregressive transformer, and the main extension is a context-and-retrieval mechanism that helps the model reuse earlier musical/chart patterns while keeping training and deployment practical through strong data and runtime optimizations.

\end{document}
